\section{Examples}

\subsection{Custom Counter}

\textbf{Requirement:} After the system enters emergency mode, at each discrete time step the counter \textit{counter\_1} is incremented by 3.

\medskip
\textbf{Entities:}
\begin{center}
\begin{tabular}{lll}
    \toprule
    \textbf{Entity} & $t_e$ & $\mu_e$ \\
    \midrule
    $\mathit{enter\_emergency}$ & EventTrigger & --- \\
    $\mathit{ev\_enter\_emergency}$ & Event & $\{\mathit{target} \mapsto id_{\mathit{enter\_emergency}},\ \mathit{type} \mapsto \mathit{generic}\}$ \\
    $\mathit{counter\_1}$ & Signal & --- \\
    \bottomrule
\end{tabular}
\end{center}

\textbf{Constraint:}
\begin{align*}
    \mathit{Always}\Bigl(\;
        \mathit{HasHappened}_\rho(\mathit{ev\_enter\_emergency},\, \mathit{Now})
        \Rightarrow \mathit{Val}_\rho(\mathit{counter\_1},\, \mathit{Now}) = \mathit{Val}_\rho(\mathit{counter\_1},\, \mathit{Prev}(\mathit{Now})) + 3
    \;\Bigr)
\end{align*}

The antecedent holds at time $t$ iff $\mathit{enter\_emergency}$ has happened at least once up to $t$.
The consequent requires $\mathit{counter\_1}$ to equal its previous value plus 3 at every such time step.

\subsection{Req 01 -- Arithmetic recurrence with condition}

\textbf{Requirement:} The signal \textit{EngagementNoSat\_u} shall equal $1.2 \times \mathit{EngagementNoSat\_u}(k{-}1)$ when \textit{Switch\_b} is true, and \textit{Lon\_u} otherwise. The initial value is 0.

\medskip
\textbf{Entities:}
\begin{center}
\begin{tabular}{ll}
    \toprule
    \textbf{Entity} & $t_e$ \\
    \midrule
    $\mathit{EngagementNoSat\_u}$ & Signal \\
    $\mathit{Switch\_b}$ & Signal \\
    $\mathit{Lon\_u}$ & Signal \\
    \bottomrule
\end{tabular}
\end{center}

\textbf{Flavour:} Discrete.

\textbf{Constraint:}
\begin{align*}
    &\mathit{Initial}\bigl(\mathit{Val}_\rho(\mathit{EngagementNoSat\_u},\, 0) = 0\bigr) \\
    \land\; &\mathit{Always}\Bigl(\\
    &\quad \bigl(\mathit{Val}_\rho(\mathit{Switch\_b},\, \mathit{Now}) = \mathit{true}\bigr) \Rightarrow {} \\
    &\qquad \mathit{Val}_\rho(\mathit{EngagementNoSat\_u},\, \mathit{Now}) = 1.2 \cdot \mathit{Val}_\rho(\mathit{EngagementNoSat\_u},\, \mathit{Prev}(\mathit{Now}))\\
    &\land\; \bigl(\mathit{Val}_\rho(\mathit{Switch\_b},\, \mathit{Now}) \neq \mathit{true}\bigr) \Rightarrow {} \\
    &\qquad \mathit{Val}_\rho(\mathit{EngagementNoSat\_u},\, \mathit{Now}) = \mathit{Val}_\rho(\mathit{Lon\_u},\, \mathit{Now})
    \;\Bigr)
\end{align*}

\subsection{Req 02 -- Store to non-volatile memory}

\textbf{Requirement:} The software shall store the maximum execution time measurement data in non-volatile memory at address \textit{MEASUREMT\_BLOCK}.

\todo{TOTO ASI UPLNE NENI TO CO CHCEME. :(}

\medskip
\textbf{Entities:}
\begin{center}
\begin{tabular}{lll}
    \toprule
    \textbf{Entity} & $t_e$ & $\mu_e$ \\
    \midrule
    $\mathit{max\_exec\_time\_data}$ & Abstract & --- \\
    $\mathit{MEASUREMT\_BLOCK}$ & Storage & $\{\mathit{non\_volatile} \mapsto \top\}$ \\
    \bottomrule
\end{tabular}
\end{center}

\textbf{Flavour:} Discrete.

\textbf{Constraint:}
\[
    \mathit{Eventually}\Bigl(
        \mathit{Val}_\rho(\mathit{MEASUREMT\_BLOCK},\, \mathit{Now}) = \mathit{Val}_\rho(\mathit{max\_exec\_time\_data},\, \mathit{Now})
    \Bigr)
\]

\subsection{Req 03 -- Ordered read then store to CPU register}

\textbf{Requirement:} The software must perform the following actions in the specified order:
1) Read the calibration constant value at the address 0xAA000018;
2) Store the calibration constant value in the DTSCON register.

\medskip
\textbf{Entities:}
\begin{center}
\begin{tabular}{lll}
    \toprule
    \textbf{Entity} & $t_e$ & $\mu_e$ \\
    \midrule
    $\mathit{calibration\_const}$ & Storage & $\{\mathit{address} \mapsto \texttt{0xAA000018}\}$ \\
    $\mathit{DTSCON}$ & Storage & $\{\mathit{register} \mapsto \top\}$ \\
    $\mathit{ev\_read\_cc}$ & Event & $\{\mathit{target} \mapsto id_{\mathit{calibration\_const}},\ \mathit{type} \mapsto \mathit{read}\}$ \\
    $\mathit{ev\_written\_dtscon}$ & Event & $\{\mathit{target} \mapsto id_{\mathit{DTSCON}},\ \mathit{type} \mapsto \mathit{written}\}$ \\
    \bottomrule
\end{tabular}
\end{center}

\textbf{Flavour:} Discrete.

\textbf{Constraint:}
\[
    \mathit{Sequence}\Bigl(
        \mathit{Happening}(\mathit{ev\_read\_cc},\, \mathit{Now}),\;
        \mathit{Happening}(\mathit{ev\_written\_dtscon},\, \mathit{Now})
    \Bigr)
\]

\todo{The Now is a little sus}

\subsection{Req 04 -- Exception vector table configuration}

\textbf{Requirement:} The software must configure IVORx registers with the addresses of Exception Vectors that compose the Exception Vector Table, according to the following table:
\begin{center}
\begin{tabular}{ll}
    \toprule
    \textbf{IVORx} & \textbf{Exception} \\
    \midrule
    IVOR0 & Critical input \\
    IVOR1 & Machine check \\
    \bottomrule
\end{tabular}
\end{center}

\medskip
\textbf{Entities:}
\begin{center}
\begin{tabular}{lll}
    \toprule
    \textbf{Entity} & $t_e$ & $\mu_e$ \\
    \midrule
    $\mathit{IVOR0}$ & Storage & $\{\mathit{register} \mapsto \top\}$ \\
    $\mathit{IVOR1}$ & Storage & $\{\mathit{register} \mapsto \top\}$ \\
    $\mathit{CriticalInput}$ & Storage & --- \\
    $\mathit{MachineCheck}$ & Storage & --- \\
    \bottomrule
\end{tabular}
\end{center}

\textbf{Flavour:} Discrete.

\textbf{Constraint:}
\[
    \mathit{Eventually}\Bigl(
        \mathit{Val}_\rho(\mathit{IVOR0},\, \mathit{Now}) = \mathit{Addr}(\mathit{CriticalInput})
        \;\land\;
        \mathit{Val}_\rho(\mathit{IVOR1},\, \mathit{Now}) = \mathit{Addr}(\mathit{MachineCheck})
    \Bigr)
\]

\subsection{Req 13 -- Bounded count of large register values}

\textbf{Requirement:} The number of values larger than 10 in the four registers A, B, C, D when they are read at the same timestep is never greater than 2.

\medskip
\textbf{Entities:}
\begin{center}
\begin{tabular}{lll}
    \toprule
    \textbf{Entity} & $t_e$ & $\mu_e$ \\
    \midrule
    $A$ & Storage & $\{\mathit{register} \mapsto \top\}$ \\
    $B$ & Storage & $\{\mathit{register} \mapsto \top\}$ \\
    $C$ & Storage & $\{\mathit{register} \mapsto \top\}$ \\
    $D$ & Storage & $\{\mathit{register} \mapsto \top\}$ \\
    $\mathit{ev\_read\_A}$ & Event & $\{\mathit{target} \mapsto id_A,\ \mathit{type} \mapsto \mathit{read}\}$ \\
    $\mathit{ev\_read\_B}$ & Event & $\{\mathit{target} \mapsto id_B,\ \mathit{type} \mapsto \mathit{read}\}$ \\
    $\mathit{ev\_read\_C}$ & Event & $\{\mathit{target} \mapsto id_C,\ \mathit{type} \mapsto \mathit{read}\}$ \\
    $\mathit{ev\_read\_D}$ & Event & $\{\mathit{target} \mapsto id_D,\ \mathit{type} \mapsto \mathit{read}\}$ \\
    \bottomrule
\end{tabular}
\end{center}

\textbf{Flavour:} Discrete.

\textbf{Constraint:}

Let $\mathit{regs} = \{A,\, B,\, C,\, D\} \in \mathcal{P}_{\text{fin}}(V_{\text{nonset}})$ be the set of register entities.
Both the antecedent and the consequent quantify over $\mathit{regs}$:

\begin{align*}
    \mathit{Always}\Bigl(
        &\mathit{ForAll}\bigl(\mathit{regs},\;
            \lambda e.\; \mathit{Happening}(e,\, \mathit{read},\, \mathit{Now})\bigr) \\
        &\Rightarrow \mathit{Cmp}\Bigl(
            \mathit{Size}\bigl(
                \mathit{Filter}\bigl(\mathit{regs},\;
                    \lambda e.\; \mathit{Val}_\rho(e,\, \mathit{Now}) > 10\bigr)
            \bigr),\; \leq,\; 2
        \Bigr)
    \Bigr)
\end{align*}

The antecedent holds at time $t$ iff all four registers are read simultaneously.
The consequent builds the set of registers whose value exceeds 10 and asserts its cardinality is at most 2.
